\subsection{CARTAN SUBALGEBRAS AND REGULAR ELEMENTS}

THEOREM 1. Let \(\mathfrak{g}\) be a Lie algebra.

\begin{enumerate}
\def\labelenumi{(\roman{enumi})}
\item
  If \(x\) is a regular element of \(\mathfrak{g}\) , \(\mathfrak{g}^0(x)\) is a Cartan subalgebra of \(\mathfrak{g}\) .
\item
  If \(\mathfrak{h}\) is a maximal nilpotent subalgebra of \(\mathfrak{g}\) , and if \(x \in \mathfrak{h}\) is regular in \(\mathfrak{g}\) , then \(\mathfrak{h} = \mathfrak{g}^0(x)\) .
\item
  If \(\mathfrak{h}\) is a Cartan subalgebra of \(\mathfrak{g}\) , then \(\dim (\mathfrak{h})\geq \mathrm{rk}(\mathfrak{g})\)
\item
  The Cartan subalgebras of \(\mathfrak{g}\) of dimension \(\mathrm{rk}(\mathfrak{g})\) are the \(\mathfrak{g}^0 (x)\) where \(x\) is a regular element.
\end{enumerate}

Let \(x\) be a regular element of \(\mathfrak{g}\) and let \(\mathfrak{h} = \mathfrak{g}^0(x)\) . Clearly \(\mathfrak{h}^0(x) = \mathfrak{h}\) . Since \(x\) is regular in \(\mathfrak{h}\) (Prop. 9), \(\operatorname{rk}(\mathfrak{h}) = \dim(\mathfrak{h})\) , so \(\mathfrak{h}\) is nilpotent. On the other hand, \(\mathfrak{h} = \mathfrak{g}^0(x) \supset \mathfrak{g}^0(\mathfrak{h}) \supset \mathfrak{h}\) , so \(\mathfrak{h} = \mathfrak{g}^0(\mathfrak{h})\) is a Cartan subalgebra of \(\mathfrak{g}\) (Prop. 4). This proves (i).

If \(\mathfrak{h}\) is a maximal nilpotent subalgebra of \(\mathfrak{g}\) , and if \(x\in \mathfrak{h}\) is regular in \(\mathfrak{g}\) , then \(\mathfrak{h}\subset \mathfrak{g}^0 (x)\) and \(\mathfrak{g}^0 (x)\) is nilpotent by (i), so \(\mathfrak{h} = \mathfrak{g}^0 (x)\) , which proves (ii).

If \(\mathfrak{h}\) is a Cartan subalgebra of \(\mathfrak{g}\) , there exists \(x \in \mathfrak{h}\) such that \(\mathfrak{h} = \mathfrak{g}^0(x)\) (Cor. 1 of Prop. 4), so \(\dim(\mathfrak{h}) \geq \operatorname{rk}(\mathfrak{g})\) , which proves (iii). If in addition \(\dim(\mathfrak{h}) = \operatorname{rk}(\mathfrak{g})\) , \(x\) is regular. Finally, if \(x'\) is regular in \(\mathfrak{g}\) , \(\mathfrak{g}^0(x')\) is a Cartan subalgebra by (i), and is obviously of dimension \(\operatorname{rk}(\mathfrak{g})\) . This proves (iv).

We shall see in §3, Th. 2 that, when \(k\) is of characteristic zero, all the Cartan subalgebras of \(\mathfrak{g}\) have dimension \(\mathrm{rk}(\mathfrak{g})\) .

COROLLARY 1. Every Lie algebra g has Cartan subalgebras, and the rank of g is the minimum dimension of a Cartan subalgebra.

COROLLARY 2. Let \(f: \mathfrak{g} \to \mathfrak{g}'\) be a surjective homomorphism of Lie algebras. If \(\mathfrak{h}'\) is a Cartan subalgebra of \(\mathfrak{g}'\) , there exists a Cartan subalgebra \(\mathfrak{h}\) of \(\mathfrak{g}\) such that \(\mathfrak{h}' = f(\mathfrak{h})\) .

Let \(\mathfrak{a} = f^{-1}(\mathfrak{h}')\) . By Cor. 1, \(\mathfrak{a}\) has a Cartan subalgebra \(\mathfrak{h}\) . By Cor. 2 of Prop. 4, \(f(\mathfrak{h}) = \mathfrak{h}'\) . We show that \(\mathfrak{h}\) is a Cartan subalgebra of \(\mathfrak{g}\) . Let \(\mathfrak{n}\) be the normalizer of \(\mathfrak{h}\) in \(\mathfrak{g}\) . It is enough to prove that \(\mathfrak{h} = \mathfrak{n}\) . If \(x \in \mathfrak{n}\) , \(f(x)\) belongs to the normalizer of \(\mathfrak{h}'\) in \(\mathfrak{g}'\) , so \(f(x) \in \mathfrak{h}'\) and \(x \in \mathfrak{a}\) ; but \(\mathfrak{h}\) is its own normalizer in \(\mathfrak{a}\) , so \(x \in \mathfrak{h}\) .

COROLLARY 3. Every Lie algebra \(\mathfrak{g}\) is the sum of its Cartan subalgebras.

The sum s of the Cartan subalgebras of g contains the set of regular elements of g (Th. 1 (i)). Since this set is dense in g for the Zariski topology, s = g.

PROPOSITION 10. Let g be a Lie algebra, a a commutative subalgebra of g and c the commutant of a in g. Assume that \(ad_{g}x\) is semi-simple for all \(x \in a\) . Then the Cartan subalgebras of c are the Cartan subalgebras of g containing a.

Let \(\mathfrak{h}\) be a Cartan subalgebra of \(\mathfrak{c}\) . Since \(\mathfrak{a}\) is contained in the centre \(\mathfrak{z}\) of \(\mathfrak{c}, \mathfrak{a} \subset \mathfrak{z} \subset \mathfrak{h}\) (Prop. 5). Let \(\mathfrak{n}\) be the normalizer of \(\mathfrak{h}\) in \(\mathfrak{g}\) . Then

\[
[ \mathfrak {a}, \mathfrak {n} ] \subset [ \mathfrak {h}, \mathfrak {n} ] \subset \mathfrak {h}.
\]

Since the \(ad_{g}x\) , \(x \in a\) , are semi-simple and commute with each other, it follows from Algebra, Chap. VIII, §5, no. 1, that there exists a vector subspace d of n stable under \(ad_{g}a\) and such that \(n = h \oplus d\) . Then \([a, d] \subset h \cap d = 0\) , so \(d \subset c\) . Thus, n is the normalizer of h in c, and hence \(n = h\) , so h is a Cartan subalgebra of g containing a.

Conversely, let \(\mathfrak{h}\) be a Cartan subalgebra of \(\mathfrak{g}\) containing \(\mathfrak{a}\) . Then \(\mathfrak{h} = \mathfrak{g}^0 (\mathfrak{h})\subset \mathfrak{g}^0 (\mathfrak{a})\) , and by hypothesis \(\mathfrak{g}_0(\mathfrak{a}) = \mathfrak{g}^0 (\mathfrak{a}) = \mathfrak{c}\) . Hence \(\mathfrak{a}\subset \mathfrak{h}\subset \mathfrak{c}\) and \(\mathfrak{h}\) is a Cartan subalgebra of \(\mathfrak{c}\) (for it is equal to its own normalizer in \(\mathfrak{g}\) , and so a fortiori in \(\mathfrak{c}\) ).

PROPOSITION 11. Let \(\mathfrak{n}\) be a nilpotent subalgebra of a Lie algebra \(\mathfrak{g}\) . There exists a Cartan subalgebra of \(\mathfrak{g}\) contained in \(\mathfrak{g}^0 (\mathfrak{n})\) .

Put \(\mathfrak{a} = \mathfrak{g}^0 (\mathfrak{n})\) . Then \(\mathfrak{n}\subset \mathfrak{a}\) since \(\mathfrak{n}\) is nilpotent. If \(x\in \mathfrak{a}\) , let \(\mathrm{P}(x)\) be the determinant of the endomorphism of \(\mathfrak{g} / \mathfrak{a}\) defined by ad \(x\) . Denote by \(\mathfrak{a}'\) the set of \(x\in \mathfrak{a}\) such that \(\mathrm{P}(x)\neq 0\) , which is an open subset of \(\mathfrak{a}\) in the Zariski topology; the relations \(x\in \mathfrak{a}'\) and \(\mathfrak{g}^0 (x)\subset \mathfrak{a}\) are equivalent. By Prop. 7 (ii) of §1, no. 2, there exists \(y\in \mathfrak{n}\) such that \(\mathfrak{g}^0 (y) = \mathfrak{a}\) , and \(y\in \mathfrak{a}'\) so \(\mathfrak{a}'\) is non-empty. Since \(\mathfrak{a}'\) is open, its intersection with the set of regular elements of \(\mathfrak{a}\) is non-empty. Let \(x\) be an element of this intersection. Then \(\mathfrak{g}^0 (x)\subset \mathfrak{a}\) and \(\mathfrak{g}^0 (x)\) is a Cartan subalgebra of \(\mathfrak{a}\) , hence is nilpotent. On the other hand, Prop. 10 (i) of §1, no. 3, shows that \(\mathfrak{g}^0 (x)\) is its own normalizer in \(\mathfrak{g}\) ; it is therefore a Cartan subalgebra of \(\mathfrak{g}\) , which completes the proof.

